\documentclass[a4paper, 12pt]{article}

\usepackage{latexsym,amsfonts,amssymb,amsthm,amsmath}
\usepackage{setspace}
\usepackage{physics}
\usepackage{graphicx}
\usepackage[utf8]{inputenc}
\usepackage[english, macedonian]{babel} % Changed from russian to macedonian
\usepackage{extsizes}
\usepackage{geometry} % Added geometry package

% Using geometry for better margin control
\geometry{
    a4paper,
    left=0.9in,
    right=0.9in,
    top=0.9in,
    bottom=0.9in,
    headheight=8pt
}

\setlength{\parindent}{0in}
\onehalfspacing

\title{
  \vspace{-1.5cm}
  \textbf{Exam problems: Hypothesis testing}
  \vspace{-1.5cm}
}
\date{}

\begin{document}

\maketitle

\section{Second midterm 2026}
\subsection*{Задача од Тестирање на хипотези – Група 1}

Ист вид на jаболка се одгледуваат во Ресен и во Тетово. Случаjно се избрани 12 стебла од Ресен и 15 стебла од Тетово, измерен е нивниот годишен принос во килограми и добиени се следниве резултати:\\


Ресен: 78, 82, 75, 80, 79, 81, 84, 77, 80, 78, 83, 79.\\
Тетово: 76, 74, 79, 77, 75, 80, 73, 78, 76, 74, 81, 75, 77, 76, 79. \\ 

Дали со ниво на значаjност 0.01 можеме да заклучиме дека очекуваниот принос на дрвjата во Ресен се разликува од очекуваниот принос на дрвjата во Тетово? Се претпоставува дека годишниот принос на jаболковите дрвjа во Ресен и Тетово се обележjа со нормална распределба и еднакви дисперзии.

\begin{itemize}
  \item [а)] (9)Задачата да се реши во тетратка!
  \item [б)] (5) Задачата да се реши во R! Кодот, заедно со излезот да се ископира подолу.
\end{itemize}

\subsection*{Задача од Тестирање на хипотези – Група 2}

(14) Еден професор предава ист предмет на две групи студенти. Едните jа следат наставата online, а другите со физичко присуство. Тоj сака да испита дали постои разлика во успехот на студентите од двете групи. Затоа случаjно избрал 11 студенти од групата со физичко присуство и 10 студенти од online групата и ги запишал поените од нивниот завршен испит:\\

Физичко присуство: 82, 75, 88, 79, 81, 85, 77, 80, 83, 78, 79\\
Online: 79, 81, 76, 84, 80, 78, 82, 77, 83, 79\\

Дали со ниво на значаjност 0.05 можеме да заклучиме дека очекуваните резултати на студентите од двете групи се разликуваат? Се претпоставува дека поените завршниот испит на двете групи се обележjа кои имаат нормална распределба и еднакви дисперзии.

\begin{itemize}
  \item [а)] (9)Задачата да се реши во тетратка!
  \item [б)] (5) Задачата да се реши во R! Кодот, заедно со излезот да се ископира подолу.
\end{itemize}


\subsection*{Задача од Тестирање на хипотези – Група 1}
Една фармацевтска компанија развила ново средство за спиење кое, како што тврдат, помага на пациентите побрзо да заспијат. Во рамки на еден клинички центар докторите одлучиле да го тестираат ова тврдење со случаен избор на 10 пациенти кои го користеле новото средство за спиење, запишувајќи ги нивните времиња на заспивање (во минути): 15, 20, 16, 20, 21, 22, 18, 17, 21 21. Вообичаено, без какво било средство за спиење, очекуваното време на заспивање на пациентите е 25 минути. Под претпоставка дека времето на заспивање има нормална распределба, со ниво на значајност од 5\% проверете дали новото средство за спиење го намалило очекуваното време на заспивање.

\subsection*{Задача од Тестирање на хипотези – Група 2}
Еден производител на батерии тврди дека стандардната девијација на животниот век на батериите е 4 месеци. Од одделот за проверка на квалитет се сомневаат дека стандардната девијација на животниот век е поголема од тоа колку што тврди производителот, што повлекува дека постои загриженост за квалитетот на батериите. За да ја проверат својата претпоставка, тие случајно избрале 12 батерии, забележувајќи го нивниот животен век: 24, 22, 28, 21, 32, 25, 30, 20, 26, 29, 33, 19. Под претпоставка дека животниот век на една батерија има нормална распределба, со ниво на значајност 5\% проверете дали стандардната девијација на животниот век на една батерија е поголема од 4 месеци.


\subsection*{Задача од ,,Параметарски тест`` – Група 1}
Тимот за аналититка во еден ланец златари коишто откупуваат злато, на случаен начин избрал 35 различни седмици и сака да спроведе анализа за масата (изразена во мг) откупено злато во текот на една седмица од страна на ланецот златари се зголемил после маркетинг кампањата што ја спровеле во текот на овие седмици. Познато е дека масата откупено злато има нормална распределба. Резултатите се следните: \\

120, 125, 130, 110, 115, 140, 135, 128, 122, 118, 132, 140, 138, 121, 133, 130, 126, 127, 129, 119, 136, 125, 123, 131, 124, 137, 139, 120, 128, 132, 125, 100, 101, 102, 103\\


Очекуваната маса откупено злато во текот на една недела пред маркетинг камањата била 95 мг. Користејќи ги прибраните податоци, со ниво на значајност 0.05, да се тестира дали очекуваната маса откупено злато во текот на една седмица се зголемила, поради маркетинг кампањата.
\begin{itemize}
  \item [а)] (7) Задачата да се реши во тетратка.
  \item [б)] (4) Задачата да се реши во R.
\end{itemize}

\subsection*{Задача од ,,Параметарски тест`` – Група 2}
Една инспекциска служба зела податоци од фабрика за производство на пијалак којшто го продава само во едно пакување, на случаен начин, по промената на машината за пакување, а поради сомнежот којашто го имала дека стандардната девијација на количеството пијалок е поголема од претходно. Од команијата тврдат дека стандардната девијација е 2. Податоците што ги добила се следните (изразени во ml): 
743, 751, 748, 757, 752, 745, 750, 755, 749, 754, 742, 756, 744, 753, 746. 
Познато е дека количеството пијалак во пакувањата има нормална распределба. Со ниво на значајност од 5\%, да се испита дали стандардната девијација на количеството пијалак во пакувањата е поголема од 2.

\begin{itemize}
  \item [а)] \textbf{(7)} Задачата да се реши во тетратка. 
  \item [б)] \textbf{(3)} Задачата да се реши во R.
\end{itemize}

\end{document}